\section{Skora's theorem: \texorpdfstring{$\protect\mathbb{R}$-trees}{Rtrees} dual to measured laminations}
\label{sec:skora}


As a final application of Theorem~\ref{thm:intersect}, we give a new
characterization of small actions of surface groups on
$\mathbb{R}$-trees, and in turn we reprove a classical theorem of
Skora~\cite{Sko90:GeometricAction}.

\subsection{Finitely generated groups acting on real trees}

We start by recalling some basic notions of actions on real trees.
An \emph{$\mathbb{R}$-tree} is a metric space $T$ such that for
all $x,y \in T$, the intersection of all connected subsets containing
both $x$ and $y$ is isometric to the interval $[0, d(x, y)]$. This
intersection is denoted
$[x, y]$, and called an \emph{arc}. If $\Gamma$ is a finitely generated group, a $\Gamma$-tree is an $\mathbb{R}$-tree together with a group action $\Gamma \acts T$ by isometries.
There are two types of isometries: those with fixed points, called \emph{elliptic}, and those without fixed points, called \emph{hyperbolic}.
Given $g\in \Gamma$, let $\| g\|\coloneqq\inf_{x \in T} d(x,gx)$, and let $T_g \coloneqq \{ x\in T : d(x,gx)=\| g\| \}$ be the \emph{characteristic set} of $g$ which is a closed, non-empty subtree of $T$ invariant under the action of $g$.
If $g$ is elliptic, then $\| g \| =0$. The converse is also true;
see~\cite[1.3]{CM87:GroupActions}.
If $g$ is hyperbolic, then $\| g \| >0$, $T_g$ is isometric to the real line and $g$ translates along $T_g$ by distance $\| g\|$. If $g$ is hyperbolic, we will call $T_g$ the \emph{axis} of $g$.

We say that $\Gamma \acts T$ is \emph{small} if the stabilizer of any (non-degenerate) arc is (virtually) cyclic.
We say that $\Gamma \acts T$ is \emph{minimal} if there is no proper non-trivial subtree of $T$ that is $\Gamma$-invariant.

An action of a finitely generated group $\Gamma$ on an $\mathbb{R}$-tree with non-trivial translation length $\| \cdot \|$ is of three types~\cite[Page~583]{CM87:GroupActions}:
\begin{enumerate}
\item Fixed end: $\| \cdot \|$ trivial on $[G,G]$.
\item Dihedral: $\| \cdot \|$ non-trivial on $[G,G]$, but $\| [g,h] \|=0$ for all $g,h$ hyperbolic in $T$.
\item Irreducible:  $\| [g,h] \| \neq 0$ for some $g,h$ hyperbolic in $T$.
\end{enumerate}

Consider two hyperbolic elements $g, h \in \Gamma$ with axes in $\mathbb{R}$-tree given respectively by $T_g,T_h$.
If $T_g \cap T_h \neq \emptyset$, we say $T_g$ and $T_h$ \emph{intersect}.

Let $\Delta \coloneqq T_g \cap T_h$. If $\Delta$ is not a singleton,
and $g$ and $h$ translate in the same direction along $\Delta$, we say
that $g$ and $h$ \emph{overlap coherently}.
If $\Delta$ is not a singleton, and $g$ and $h$ translate in opposite
directions along $\Delta$, we say that $g$ and $h$ \emph{overlap incoherently}.
Following~\cite[3.2]{CM87:GroupActions}, we say that $g,h$ form a \emph{good pair} if 
\begin{itemize} 
\item they are both hyperbolic;
\item their axes in the tree intersect and their overlap is not a singleton;
\item they overlap coherently; and
\item the length of $\Delta$ is strictly less than $\min \{ \| g \|, \| h \| \}$.
\end{itemize}

\begin{remark}[{\cite[Remark~3.3]{CM87:GroupActions}}]
If $g,h$ are hyperbolic and $T_g,T_h$ meets in an arc of positive but finite length, then there exist integers $m>0$ and $n$ so that $g^m$ and $h^n$
for a good pair of isometries, where the sign of $n$ is chosen so that the overlap is coherent, and the magnitudes of $m,n$ are chosen
to ensure $m \| g \|$ and $n \|h\|$ are both greater than the length  of $T_g \cap T_h$.
\label{rmk:powers_goodpair}
\end{remark}

\begin{lemma}
Let $\Gamma$ be a finitely generated group. $\Gamma$ on $T$ is irreducible if and only if there exists at least one good pair.
\label{lem:char_irred_good}
\end{lemma}
\begin{proof}
By~\cite[Theorem~2.7]{CM87:GroupActions} $\Gamma$ on $T$ is irreducible if and only if there exist a pair of hyperbolic elements whose axes have bounded non-empty intersection. By~\cite[Theorem~3.3]{CM87:GroupActions} this implies the existence of a good pair. On the other hand, one of the requirements of a good pair is that both elements are hyperbolic and their axes have bounded overlap, which again by~\cite[Theorem~2.7]{CM87:GroupActions} implies $\Gamma$ on $T$ is irreducible. 
\DPTtodo[inline]{Why does this imply $\norm{[g,h]} \ne 0$?}
\DMGtodo{added the last line to be more explicit}
\end{proof}


\begin{lemma}
  For any action of a finitely-generated group on an $\bbR$-tree, if $g,h$ form a good pair, then $gh^{-1}$ and $h^{-1}g$ are hyperbolic isometries of $T$ and $T_{gh^{-1}}, T_{h^{-1}g}$ do not intersect.
    \label{lem:oppgoodpair}
\end{lemma}
\begin{proof}
   This is contained in the proof of~\cite[Lemma~3.4]{CM87:GroupActions}.
\end{proof}


Another way to characterize irreducibility is in terms of the existence of a non-intersecting pair of axes.
\begin{lemma}
A finitely generated group $\Gamma$  acting on a real tree $T$ is irreducible if and only if there exist hyperbolic elements $g,h$ in $T$ so that $T_g \cap T_h = \emptyset$.
\label{lem:char_irred_noninter}
\end{lemma}
\begin{proof}
Follows by combining Lemmas~\ref{lem:char_irred_good} and \ref{lem:oppgoodpair}.
\end{proof}



In general, if the translation lengths of two $G \acts T$ actions are equal, tha does not necessarily imply the trees are equivariantly isometric, but it is true if the action is irreducible by~\cite[Theorem~3.7]{CM87:GroupActions}.
In fact, if the length is non-zero, and the action is assumed irreducible, by \cite[Proposition~3.1]{CM87:GroupActions} there is a unique minimal invariant subtree (i.e., a subtree invariant under $\Gamma$ that does not contain any proper non-degenerate invariant subtree). The length function for the restriction to this minimal invariant subtree is equal to the length function of $G  \acts T$. Hence, without loss of generality, one can always restrict to the action on the minimal invariant subtree without changing the length functional. In what follows. we will always consider minimal actions.

We will consider the following two handy results about the relative position of axes in a real tree, which can be found in~\cite[1.5, 1.8]{CM87:GroupActions}.

\begin{lemma}
\begin{enumerate}
\item (Intersecting axes)
If $g,h \in \Gamma$ are hyperbolic, then
$T_g \cap T_h \neq \emptyset$ if and only if $\| g \| + \| h \| =\max \{ \| gh \| , \|gh^{-1}\| \}$.
Moreover, if $g,h$ overlap coherently, then 
$\| g \| + \| h \| = \| gh \| >\|gh^{-1}\|$,
and if $g,h$ overlap incoherently, then 
$\| g \| + \| h \| = \| gh^{-1} \| >\|gh\|$.
\item (Parallel axes) If $g,h$ are any two isometries satisfying $T_g \cap T_h =\emptyset$ or $g,h$ are hyperbolic and $T_g \cap T_h = \{ * \}$, then
\[
\|gh\| = \| gh^{-1}\| = \|  g \| + \|h \| + 2d(T_g,T_h).
\]
\label{lem:foundation_trees}
\end{enumerate}
\end{lemma}

We will also use the following, found in~\cite[Lemma~6.2]{CM87:GroupActions}

\begin{lemma}
    If $g,h$ are hyperbolic and satisfy $\|g\| + \|h\| = \|gh\|$, then for every positive integers $m,n$, we have
    \[
    \|g^m h^n\| = m \| g \| + n \|h \|.
    \]
    \label{lem:powers_intersect}
\end{lemma}

\subsection{Surface group actions on real trees}

In this section, we assume $\Gamma = \pi_1(S)$.

It is standard that a measured lamination $\lambda$ defines a (dual) minimal irreducible $\Gamma$-action on $\mathbb{R}$-tree (see, for example, \cite[Chapter~11]{Kap09:HyperbolicManifold}).
It follows from the construction of $T_\lambda$ that the translation length on $T$ of $g \in \pi_1(S)$, $\| g \|$, is equal to $i(\lambda,[g])$, where $[g]$ denotes the free homotopy class of a closed curve associated to $g$.
Given a minimal action of a surface group $\Gamma$ on an $\mathbb{R}$-tree, if there is a measured lamination $\lambda$ and $\Gamma$-equivariant isometry $\phi \colon T_\lambda \to T$, we say that $T$ is \emph{dual to the measured lamination} $\lambda$.

If $\Sigma$ is a hyperbolic structure on $S$, let $\ora{g}, \ora{h}$ be the hyperbolic axes corresponding to these elements.
If two axes $\ora{g},\ora{h}$ with respect to $\Sigma$ intersect transversely in $\tilde{\Sigma}$, the axes with respect to any other hyperbolic structure also intersect. 

Morgan-Otal~\cite{MO93:RelativeGrowths} first showed that a minimal action of a surface group is dual to a measured lamination if and only if 
\begin{enumerate}
    \item it is small,
    \item it \emph{preserves axis intersection}: for any pair $g,h \in \Gamma$ acting as hyperbolic isometries on $T$, we have $T_g \cap T_h \neq \emptyset$ whenever $\ora{g}, \ora{h}$ cross.
\end{enumerate}

Morgan-Otal asked in~\cite{MO93:RelativeGrowths}: if a minimal action preserves axis intersection, does it follow it is small?
Without further assumptions, the axis intersection being preserved alone does not imply in general that the action is small, as the example of a dihedral action of a surface group acting on $\mathbb{R}$ shows (this was noted by Skora in~\cite[Section~4]{Sko90:GeometricAction}).
However, in~\cite{Sko90:GeometricAction}, Skora showed that a minimal action $\Gamma \acts T$ of a surface group $\Gamma$ is dual to a measured lamination if and only if it is irreducible and preserves axis intersection.
In~\cite{Sk96:Splittings},  Skora showed that $\Gamma \acts T$ is dual to a measured lamination if and only if the action is small. In fact, Skora showed this in the more general setting where $S$ has cusps provided that then parabolic elements in $\Gamma$ are elliptic in $T$.
It follows from this that the action being small is equivalent to it being irreducible and preserving axis intersection.

The main result of this section will be the following.

\begin{theorem}
Let $\pi_1(S) \acts T$ on a $\mathbb{R}$-tree $T$, irreducible and that preserves axis intersection. Then it follows that $\| \cdot \|$ satisfies max-smoothing.
\label{prop:small}
\end{theorem}

Even though we will not needed it, as a warm-up, it is worth noticing that a reducible action does not satisfy smoothing.

\begin{lemma}
Let $\Gamma=\pi_1(S)$ be a surface group acting on a real tree $T$. Let $\| \cdot \|$ denote the translation length of this action.
If $\| \cdot \| \colon \mathcal{C}(S) \to \mathbb{R}$ satisfies smoothing, then $(\Gamma, T)$ is irreducible.
\end{lemma}
\begin{proof}
We start with a reducible action and show it cannot satisfy smoothing.
If the action is reducible, every pair $(g,h)$ of elements acting hyperbolically on $T$ has axes overlapping in a ray. By~\cite[Proposition~1.6~(3)]{Pa89:Paulin1989TheGT}, either $g,h$ overlap coherently, in which case
\begin{equation}
\|gh \|=\| g \| + \| h \| \mbox{ and } \|gh^{-1} \| = |\|g \|- \| h \| |,
\label{eq:reducible1}
\end{equation}
or $g,h$ overlap incoherently, in which case
\begin{equation}
\|gh^{-1} \| =\| g\| + \| h \| \mbox{ and } \| gh \| = | \| g \| - \| h \||.
\label{eq:reducible2}
\end{equation}

Take a pair $a,x \in \Gamma$, so that $(a,x)$ $R$-cross.
WLOG, assume $\| a \|<\| x \|$.

\begin{claim}
Without loss of generality, we can furthermore assume that $a$ and $x$ overlap coherently.
\end{claim}
\begin{proof}[Proof of Claim]
If $a,x$ overlap incoherently, then note $(xa^{-1}, ax)$ also $R$-cross, by Lemma~\ref{lem:righthandedfoundation}. 
Then $\| ax \| =\| x\|- \|a \|$ and $\| xa^{-1} \| =\| a \| + \| x \|$. Since $\| (xa^{-1})(ax) \| =\| x^2 \|$ is equal to $\| ax \| + \| xa^{-1} \|= 2\| x \|$, $ax$ and $xa^{-1}$ overlap coherently, and they do it in a subray of the overlap between $a$ and $x$.
\end{proof}
If $(a,x)$ R-cross, we use Equation~\ref{eq:cross-split} in Proposition~\ref{prop:join-split-geom}, with $m=-n$, $a=b$, and $n$ large enough, and choosing some admissible $r$ for the inequalities there, by smoothing,
\begin{equation}
 \| ax^nax^{-n} \| \geq \| ax^{n+r} \| + \| ax^{-n-r} \|.
 \label{eq:unoriented1}
\end{equation}
On the other hand, by a similar argument as in the proof of the above Claim, $ax^n, ax^{-n}$ overlap incoherently, so, we have that the left hand side is
\begin{equation}
\|ax^nax^{-n} \| = \| ax^n \| - \|ax^{-n} \| = \| a \| + n\| x \| - (n\| x \|-\|a\|) = 2\|a\|
\label{eq:lhs1}
\end{equation}
where we used that Equation~\ref{eq:reducible2} applies to the pair ($ax^n,ax^{-n}$) overlapping incoherently in the first equality, and Equation~\ref{eq:reducible1} for the coherent pair $(a,x^n)$ and Equation~\ref{eq:reducible2} for the incoherent pair $(a,x^{-n})$ in the second equality.
A similar computation, yields the following result for the right hand side:
\begin{equation}
\|ax^{n+r}\|=\|a\|+(n+r)\|x\| + (n+r)\|x\|-\|a\| = 2(n+r)\|x\|.
\label{eq:rhs1}
\end{equation}
Combining Equation~\ref{eq:unoriented1} with Equations~\ref{eq:lhs1} and~\ref{eq:rhs1} we get $\|a \| \geq (n+r)\|x\|$, which contradicts our assumption.
\end{proof}

\begin{remark}
Note that by Equations~\ref{eq:reducible1}~and~\ref{eq:reducible2} a reducible action always satisfies disconnected smoothing, but it never satisfies connected smoothing. Since a reducible action cannot be dual to a measured lamination, this shows once again that one cannot strengthen Theorem~\ref{thm:intersect} by assuming only disconnected smoothing. Another example in the setting of the linear extension of extremal length can be seen in Example~\ref{ex:linEL}.
\label{rmk:reducible_action}
\end{remark}




Let us now address the proof of Theorem~\ref{prop:small}.

\begin{proof}[Proof of Theorem~\ref{prop:small}]
We assume $(\Gamma, T)$ is irreducible and preserves axis intersection, and will prove $\| \cdot \|$ satisfies max-smoothing. 
Assume that for any hyperbolic pair $g,h \in \pi_1(S)$ on $T$ we have $T_g \cap T_h \neq \emptyset$ whenever $\ora{g}$ and $\ora{h}$ cross, where $T_g$ denotes the axis of $g$ in the tree.
We will verify that $\| \cdot \|$, the translation length on $T$, satisfies the max-smoothing condition, and thus it comes from a measured lamination.
There are several cases to consider, depending on the relative position of the axes of the elements in the tree. For each case, different inequalities between $\| g\| + \| h \|$, $\| gh \|$ and $\| gh^{-1} \|$ are satisfied, and were analyzed by Paulin~\cite{Pa89:Paulin1989TheGT}. The strategy of the proof will be to compare the intersection of axes in the surface versus those in the tree, for each one of those cases.
We will see that max-smoothing is not satisfied in all cases. In such cases, we will find, using Lemma~\ref{lem:oppgoodpair}, pairs of elements whose axes do not intersect in the tree, but they do in the surface, which will allow us to rule them out, as that would contradict the assumption of preservation of axis intersection.

In order to prove Theorem~\ref{prop:small}, we consider several cases.

\begin{enumerate}[wide, labelindent=0pt]
\item Case 1 (hyperbolic): we suppose that for a pair of elements $g,h \in \pi_1(X)$, all $g,h$ act hyperbolically on $T$.
\begin{enumerate}
    \item Case 1a. First, suppose that $\ora{g}$ and $\ora{h}$ cross.
The two possible smoothings are then represented by $gh$ and $gh^{-1}$, so max-smoothing can be phrased as

\[
\|g \| + \| h \| = \max \{ \| gh \|, \| gh^{-1} \| \}.
\]

This relation is true by~\cite[1.8]{CM87:GroupActions}, since $T_g \cap T_h \neq \emptyset$.
\item Case 1b. Suppose that $\ora{g}$ and $\ora{h}$ are parallel.
Max-smoothing amounts to

\[
\| gh \| = \max \{ \| g \| +\| h \| , \| gh^{-1} \| \}.
\]

 We have that $\ora{gh}$ and $\ora{hg}$ cross. Thus, since the action preserves axis intersection, we know that $T_{gh} \cap T_{hg} \neq \emptyset$.
There are three subcases to consider.
\begin{enumerate}
    \item $T_g \cap T_h = \emptyset, \{ p \}$. Then $\|gh\|=\|gh^{-1}\| \geq \|g\|+\|h\|$ by the result~\cite[1.5]{CM87:GroupActions}, and the condition is satisfied.
    \item $T_g$ and $T_h$ intersect in a non-degenerate interval with coherent orientations.
    Then, 
    $\|gh\|=\|g\|+\|h\|>\|gh^{-1}\|$, and the condition is satisfied by the result~\cite[1.8]{CM87:GroupActions}
    \item  $T_g$ and $T_h$ intersect in a non-degenerate interval with opposite orientations. Then, we would have
    $\|gh^{-1}\|=\|g\|+\|h\|>\|gh\|$, so the max-smoothing condition would be violated.
    We show this cannot happen. 
    Within this case, there are two subcases to consider. Let $\Delta$ denote the length of the overlap between $T_g$ and $T_h$.
    \begin{enumerate}
        \item If $\Delta<\min \{ \| g \|, \| h \| \}$, then $g, h^{-1}$ form a good pair, and thus by Lemma~\ref{lem:oppgoodpair} we have $T_{gh} \cap T_{hg} = \emptyset$, a contradiction.
        \item If $\Delta \geq \min \{ \| g \|, \| h \| \}$,
        then, according to Remark~\ref{rmk:powers_goodpair} we take $m>0,n>0$ so that $g^m$ and $h^{-n}$ form a good pair, and hence, Lemma~\ref{lem:oppgoodpair}, we have $T_{g^mh^n}$ and $T_{h^ng^m}$ do not intersect. Since $\|g \| + \|h \| = \|gh\|$, by Lemma~\ref{lem:powers_intersect}, we have
        \[
          \|g^m h^n \| = m\|g\| + n\|h\|
        \]
        which, by Lemma~\ref{lem:foundation_trees} (i), implies $T_{g^mh^n}$ and $T_{h^ng^m}$  intersect, a contradiction.
\end{enumerate}
    
    \end{enumerate}
        \end{enumerate}
\item Case 2: $g$ or $h$ are elliptic.
It follows from a similar analysis comparison of translation lengths in the tree. The different cases when, say, $g$ is elliptic and $h$ is hyperbolic, have been analyzed in~\cite[Proposition~1.7]{Pa89:Paulin1989TheGT}. 

\begin{enumerate}
\item Suppose $T_g \cap T_h = \emptyset$. 
\begin{itemize}[wide, labelindent=0pt] 
\item In the disconnected smoothing case, we show that this case, in fact, never occurs. Indeed, if $T_g \cap T_h = \emptyset$, it follows by~\cite[Lemma~3.2]{Chi01:LambdaTrees} that $gh$ is hyperbolic, and in fact,
 \[
 \| g h \| = \| gh^{-1} \| = \| g \| + \| h \| + 2d(T_g, T_h)
 \]
 and if $\Delta_1$ denotes the length of $T_{gh} \cap T_h$ and $\Delta_2$ denotes the length of $T_{gh} \cap T_g$, then $\Delta_1=\|h\|$ and $\Delta_2=\|g\|$.
 Thus, by the equation above, we have that $\|gh \| > \|h\|$.
 Then we are in~~\cite[Proposition~1.7]{Pa89:Paulin1989TheGT}, case (2). We reproduce the more complete configuration of axes in Figure~\ref{fig:tree}.
 
      \begin{figure}
\centering{
\svgfignormal{tree}
    \caption{Configuration of axes when $\protect\ora{g}$ and $\protect\ora{h}$ are not intersecting}\label{fig:tree}}
\end{figure}

We observe that $g h$ and $(h g)^{-1}$ overlap in $\Delta=d(z,y)$ so they form a good pair. Thus,  by Lemma~\ref{lem:oppgoodpair}, $T_{g h^2 g}$ and $T_{h g^2 h}$ do not intersect.
On the other hand, since $\ora{g},\ora{h}$ cross, then $\ora{gh}$ and $\ora{hg}$ are parallel, by Lemma~\ref{lem:righthandedfoundation}, part 1, and thus $\ora{gh^2g}$ and $\ora{hg^2h}$ intersect by Lemma~\ref{lem:righthandedfoundation}, part 2, so we have a contradiction. This finishes the proof for the disconnected smoothing case.
\item  In the connected smoothing case, the formula above says $\| g h \| = \| gh^{-1} \|$, so max-smoothing is satisfied.
\end{itemize}
\item  Suppose $T_g \cap T_h \neq \emptyset$. 
\begin{itemize}[wide, labelindent=0pt]
\item Suppose $T_h$ and $T_{ghg^{-1}}$ are equal. Max smoothing would fail (both connected and disconnected) if $g$ acted as a reflection about the only fixed-point of $g$ on $T_h$, since then $\| gh \|=\| g^{-1}h \|=0$. 
But in that case, $gh$ (and $hg$) are elliptic, and $gh^2g$ is hyperbolic (with hyperbolic axis equal to $\ora{h}$).
In the connected smoothing case, we can apply Case 3 (below) to $gh, hg$ in the \emph{disconnected} smoothing case, since $\ora{gh},\ora{hg}$ intersect, ruling this case out.
In the disconnected smoothing case, we can similarly apply Case 3 (below) to $g, hg$ in the \emph{disconnected} smoothing case
since $\ora{g},\ora{hg}$ intersect, and rule this case out.
\item Suppose $T_h$ and $T_{ghg^{-1}}$ intersect in a compact segment.

For disconnected max-smoothing we have that $\ora{g}$ and $\ora{h}$ cross, and we want to see $\| h \| = \max \{ \|gh\| , \|gh^{-1}\| \}$.
By ~\cite[Proposition~1.7]{Pa89:Paulin1989TheGT}, we must have that either $T_{g h g^{-1}}$ and $T_{h}$ intersect coherently, or $T_{g h g^{-1}}$ and $T_{h^{-1}}$ intersect coherently. In the first case, $\| gh \| = \| h \|$, and $\| h \| = \| h^{-1} \|\geq \| gh^{-1} \|$, by~\cite[Proposition~1.7]{Pa89:Paulin1989TheGT}, case (1), (ii), and in the second case, $\| gh^{-1} \| = \| h^{-1} \|=\| h \|$, and $\| h \| \geq \| gh \|$. Thus, disconnected max-smoothing is satisfied in any case.

In the connected smoothing case, $\ora{gh}$ and $\ora{hg}$ intersect, so~$\| gh \| = \max \{ \|h\| , \|gh^{-1}\| \}$.
Again, by ~\cite[Proposition~1.7]{Pa89:Paulin1989TheGT}, we must have that either $T_{g h g^{-1}}$ and $T_{h}$ intersect coherently, or $T_{g h g^{-1}}$ and $T_{h^{-1}}$ intersect coherently. 

In the first case, $\| gh \| = \| h \|$, and $\| h \| = \| h^{-1} \|\geq \| gh^{-1} \|$, so max-smoothing is satisfied. 
In the second case, $\| gh^{-1} \| = \| h^{-1} \|=\| h \|$, and $\| h \| \geq \| gh \|$, which would contradict max-smoothing.
However, we can rule this out. First, suppose that $\| h \| > \Delta$, where $\Delta$ is the overlap between $T_{h}$ and $T_{g h g^{-1}}$. Building from Fig.7~\cite[Proposition~1.7]{Pa89:Paulin1989TheGT}, it follows that $T_{gh}$ and $T_{hg}$ are disjoint, as Figure~\ref{fig:oppositehypell} illustrates. 
It remains to justify that $g(x)=y$ and $x=g(y)$. We start by proving the following, which is implied by  Fig.7~\cite[Proposition~1.7]{Pa89:Paulin1989TheGT} but not proven.

\begin{claim}Let $z$ be the unique fixed point of $h$ in the interval $[x,y]$. Then $z$ is the middle point between $x$ and $y$, and $g(x)=y$, $g(y)=x$.\end{claim}
\begin{proof}[Proof of Claim]
Following Paulin's notation, let $[z, +\infty)_{h}$ (resp.\ $(-\infty,
z]_{h}$)  be the open half-line of $\ora{h}$, with endpoint $z$, that
does not contain $x$ (resp.\ $y$) but
contains $y$ (resp.\ $x$). In the same way, let $[z,
+\infty)_{ghg^{-1}}$ (resp.\ $(-\infty, z]_{ghg^{-1}}$ be the
open half-line of $\ora{ghg^{-1}}$ with endpoint $z$, that does not contain $x$ (resp.\ $y$) but
contains $y$ (resp.\ $x$). Since $g$ sends $[z, +\infty)_{h}$ to $(-\infty, z]_{ghg^{-1}}$ and sends $(-\infty, z]_{ghg^{-1}}$ to $[z, +\infty)_{ghg^{-1}}$, it follows that $[z,y]$ must be sent to $[z,x]$, and $[z,y]$ to $[z,x]$. Thus, the two segments must have equal length, half the length of the overlap $\Delta$, so $z$ must be the middle point of $[x,y]$ and $g(x)=y$, $g(y)=x$. 
\end{proof}

If $\| h \| \leq \Delta$, then we take $n>0$ large enough so that $\| h^n \| > \Delta$.
By the same analysis as before, it follows that $T_{h^ng}$ and $T_{gh^n}$ are disjoint. On the other hand, an application of Lemma~\ref{lem:righthandedfoundation}, shows that $\ora{h^n g}$ and $\ora{gh^n}$ intersect for all $n>0$, which is a contradiction.
      \begin{figure}
\centering{
\resizebox{80mm}{!}{\svgfig{oppositehypell}}
    \caption{Configuration of axes when $\protect\ora{ghg^{-1}}$ and $\protect\ora{h}$ overlap incoherently with overlap less than $\| h \|$.}\label{fig:oppositehypell}}
\end{figure}
\item Suppose $T_h$ and $T_{ghg^{-1}}$ intersect in a ray $R_{\infty}$ with endpoint $x$.
This is similar to the preceeding subcase.
\begin{comment} If both $h$ and $ghg^{-1}$ translate in the same direction along the ray $R_{\infty}$, then $\| gh \| = \| h \|$ and max-smoothing is satisfied in (both connected and disconnected). If they translate in opposite directions, and $\gamma x \notin R_{\infty}$ \end{comment}

\end{itemize}
\end{enumerate}
\item Case 3. Suppose $g$ and $h$ are elliptic. We follow~\cite[Proposition~1.8]{Pa89:Paulin1989TheGT}.
If $gh$ is elliptic, we are done in both connected and disconnected smoothing. \footnote{In fact, this case does not arise, since it would contradict irreducibility, by~\cite[Lemma~10.4]{Kap09:HyperbolicManifold}}
If $gh$ is hyperbolic, then the diagram in the proof of~\cite[Proposition~1.8]{Pa89:Paulin1989TheGT} shows that
$T_{gh}$ and $T_{hg}$ overlap in $[z,z']$ incoherently. Indeed, the overlap can only be in $[z,z']$, since if $ w \in T_{hg} \cap [z',gz] \neq \emptyset$, then $hg(w)=w'$ in Figure~\ref{fig:ell}. But on the other hand we should have $hg(w) \in [hg z', hg z]$, since $w \in [\delta(z'), \delta(z)]$, because $\delta(z')=z'$.
Given that $\| gh \| = 2d(z,z')$, this implies that $gh$ and $(hg)^{-1}$ form a good pair, and by Lemma~\ref{lem:oppgoodpair},
$T_{gh^2g}$ and $T_{hg^2h}$ do not intersect.
In the disconnected smoothing case, when $\ora{g}$ and $\ora{h}$ intersect, an application of Lemma~\ref{lem:righthandedfoundation} part 2, and then part 1, shows that $\ora{gh^2g}$ and $\ora{hg^2 h}$ intersect, which contradicts preservation of axis intersection.
In the connected smoothing case, we get $\| gh \| = \| gh^{-1} \|$, so max-smoothing is verified.

      \begin{figure}
\centering{
\resizebox{60mm}{!}{\svgfig{ell}}
    \caption{Configuration of axes when $g$ and $h$ are elliptic}\label{fig:ell}}
\end{figure}


\end{enumerate}
\end{proof}

Now we prove Theorem~\ref{mainthm:skora}.

\begin{proof}[Proof of Theorem~\ref{mainthm:skora}]
(A) is equivalent to (C) is proven by the first part of Theorem~\ref{mainthm:apps}.
(B) implies (C) is Theorem~\ref{prop:small}.
(C) implies (B) follows by work of Morgan-Otal~\cite{MO93:RelativeGrowths}.
\end{proof}

\begin{remark}

Cases 2 and 3 in the above proof seem necessary, in the sense that, in principle, the proof cannot be simplified by just proving Case 1 and approximating by minimal and free actions.
Indeed, note that minimal free actions are small and Culler-Morgan~\cite{CM87:GroupActions} proved that the subspace of minimal small actions of any finitely generated group is closed in the space of all actions.
It follows that the closure of minimal free actions is contained in the space of minimal small actions. To simplify the proof one would have to prove this closure is \emph{the same as} minimal small actions, a fact we know is true a posteriori, for example, because maximal measured laminations, (whose actions are dual to free minimal actions on real trees)~\cite{FLP79:ThurstonWorkSurfaces} are dense in the space of all measured laminations.
On the other hand, note that the analogous result is not true in the setting of free groups: by a result of Bestvina-Feighn~\cite{BF11:OuterLimits}
the closure of the space of minimal, free, simplicial actions is the space of so-called ``very small'' actions, which is a proper subset of small actions.
In the setting of surface groups, the space of very small and small actions coincide, so the density result of free minimal actions holds.
\end{remark}

\subsection{Inequalities between intersection numbers}
\label{subsec:ineqintersection}

Finally, we give an application of our main result to domination of spectra of metric lengths on a surface. Let $\mu_1,\mu_2 \in \Curr(S)$. Say $\mu_1 \leq \mu_2$, or $\mu_2$ \emph{dominates} $\mu_1$ if $i(\mu_1,C) \leq i(\mu_2,C)$ for all $C \in \mathcal{C}(\Sigma)$.
We can use these extension results to study what intersection numbers dominating one curve over another imply about Liouville currents.

For example, we answer positively the following question at the end of \cite{NC01:ShortestGeodesics}.
In fact, we have something much stronger

\begin{corollary}
If $C_2$ dominates $C_1$, then
    $\ell_g(C_1) \leq \ell_g(C_2)$ for all length metrics $g$ on $S$.
    \label{cor:domination}
\end{corollary}

\begin{proof}
First note that any length metric satisfies smoothing.
This argument is outlined essentially in \cite[Section~4.1]{MGT21:Smoothings}, but we reproduce it here.
 Take a
curve~$C$ with an
essential self-intersection~$x$, and a representative~$\gamma$. Then $\gamma$ has
a self-intersection point~$x'$ of the homotopy type of~$x$. If we
consider the
curve representative~$\gamma' \in C'$ obtained by smoothing at~$x'$, then,
since $\ell_g$ is an infimum, we have
\[
  \ell_g(C')\le \ell_g(C') = \ell_g(\gamma) = \ell_g(C),
\]
as desired. $\ell_g$ satisfies stability by \cite[Lemma~4.2]{MGT21:Smoothings}. Additivity is satisfied by definition of $\ell_g$ on weighted multi-curves by extending linearly.
It is clear that $\ell_g$ is a symmetric curve functional, since it is naturally defined on unoriented curves.
Thus, by Theorem~\ref{thm:intersect}, there exists a geodesic current $\mu_g$ so that $\ell_g(C)=i(\mu_g,C)$ for all closed curves $C$. Approximating the current $\mu_g$ by weighted multi-curves and using continuity of intersection number on currents the result follows.
\end{proof}


%%% Local Variables:
%%% mode: latex
%%% TeX-master: "Intersections"
%%% End:
